\documentclass[10pt,a4paper]{amsart}
\usepackage[margin=19mm]{geometry}
\usepackage{amsmath,amssymb,mathtools}
\usepackage[scr=rsfs,cal=euler]{mathalfa}
\usepackage{xfrac}
\usepackage[unicode,hidelinks,bookmarks=false]{hyperref}
\newcommand{\ie}{i.e.\ }
\newcommand{\cf}{cf.\ }
\newcommand{\resp}{respectively\ }
\newcommand{\Subsection}{\subsection}
\providecommand{\nobreakdash}{\nobreak\hbox{-}}
\setlength{\emergencystretch}{2em}
\multlinegap0pt
\usepackage{bm}
\usepackage{bbm}
\usepackage{dsfont}
\usepackage{xspace}
\usepackage{cancel}
\usepackage{mathtools}
\usepackage[shortlabels]{enumitem}
\usepackage{amsthm}
\makeatletter
\@ifundefined{Definition}{\newtheorem{Definition}{Definition}}{}
\@ifundefined{Lemma}{\newtheorem{Lemma}{Lemma}}{}
\@ifundefined{Proposition}{\newtheorem{Proposition}{Proposition}}{}
\@ifundefined{Theorem}{\newtheorem{Theorem}{Theorem}}{}
\makeatother
\title[Center of distances of ultrametric spaces generated by label]{Center of distances of ultrametric spaces generated by labeled trees}
\author{Oleksiy Dovgoshey, Olga Rovenska}
\date{Source: arXiv:2601.13363v2 (CC BY 4.0)}
\begin{document}
\maketitle

\noindent\emph{Source and licence.} Center of distances of ultrametric spaces generated by labeled trees. Version arXiv:2601.13363v2 (CC BY 4.0), distributed under the Creative Commons Attribution 4.0 International licence, as recorded in the arXiv licence metadata for this exact version and shown on the abstract page https://arxiv.org/abs/2601.13363v2. Full rights notice: licenses/arxiv-2601.13363-CC-BY-4.0.txt.\par\smallskip

\noindent\textit{Editorial scope.} Counted unit: the Theorem whose source label is \texttt{efhik} in the article (source0.tex lines 1252--1266, source label \texttt{efhik}), delivered together with the article's own complete original proof of it (source0.tex lines 1269--1356, one \texttt{proof} environment of 88 lines). No mathematics is written, repaired or re-derived: statement and proof are the article's own text line for line. Required context retained uncounted: lokias (lines 264--266); nnhh (lines 312--318); scfr (lines 434--438); ref (lines 466--474); dggkq (lines 806--827); sat (lines 958--960); shak1 (lines 968--976); scgkk (lines 1066--1074); dhiu (lines 1078--1080). Every definition, display and label of those lines is retained unchanged. Omitted from this package and not redistributed: 1--263, 267--311, 319--433, 439--465, 475--805, 828--957, 961--967, 977--1065, 1075--1077, 1081--1251, 1267--1268, 1357--1853. Nothing inside a retained excerpt is omitted or altered: every retained body is delivered exactly as the article's lines are written, so no omission receipt of any kind accompanies this package. Citations: the retained excerpts carry no citation command at all, so no bibliography entry of the article is transcribed and no reference is redistributed. Reference adaptation: none was needed and none was made; every reference inside the delivered unit resolves inside it, so no unresolved reference or citation is produced. Printed slips kept literal: no printed word, symbol, hypothesis or number of the counted statement or of its proof was altered. Layout and class: the article's own class is \texttt{article} with packages including color, cite, amsfonts, amssymb, amsmath, mathrsfs, textcomp, amsbsy, amsthm, enumitem, lineno, graphicx, fontenc, tikz; the wrapper is the lane's pinned \texttt{amsart} editorial wrapper at 10pt on A4, which restates the theorem-like environments the retained text needs and adds the article's own macro definitions that the retained text needs (none), with extra wrapper packages (bm, bbm, dsfont, xspace, cancel, mathtools, enumitem). The delivered numbering is the unit's own. Assets: no retained span contains a figure environment, a graphics-inclusion command, a PDF-page inclusion, a table or a TikZ picture; no image file is copied into this package, the package declares no asset dependency and no image file is redistributed with it. Licence: version arXiv:2601.13363v2 of this article is distributed under the Creative Commons Attribution 4.0 International licence, as recorded in the arXiv licence metadata for this exact version and shown on the abstract page https://arxiv.org/abs/2601.13363v2. A full rights notice is delivered at licenses/arxiv-2601.13363-CC-BY-4.0.txt. Characters: every retained excerpt is pure ASCII, so the delivered engine, XeLaTeX with its default Latin Modern text font, reports no missing character.
\begin{equation}\label{lokias}
    \operatorname{diam} A := \sup\{ d(x,y) : x,y \in A \}
\end{equation}

\begin{Lemma}\label{nnhh}
Let $A$ be a non-empty subset of an ultrametric space $(X,d)$. Then the equality
\begin{equation}\label{ssqq}
\operatorname{diam} A=\sup\{d(p,a): a\in A\}
\end{equation}
holds for each $p\in A$.
\end{Lemma}

\begin{Definition}\label{scfr} Let $G$ be a graph and let $k \geq 2$ be a cardinal number. The graph $G$
is complete $k$-partite if the vertex set $V(G)$ can be partitioned into $k$ non-empty
disjoint parts, in such a way that no edge has both ends in the same part and
any two vertices in different parts are adjacent.
\end{Definition}

\begin{Theorem}\label{ref} Let $(X,d)$ be an ultrametric space with $|X| \geq 2$. Then the following
statements are equivalent:

\begin{enumerate}[label=\textit{(\roman*)}, left=0pt]
\item The diametrical graph $G_{X}$ is non-empty.
\item The diametrical graph $G_{X}$ is complete multipartite.
\end{enumerate}

\end{Theorem}

\begin{Theorem}\label{dggkq}
Let $(X,d) \in {\bf UT}$ have at least two points, let $D(X)$ be the distance set of $(X,d)$, and let $C(X)$ be the center of distances of $(X,d)$.
Then
 we have either
\begin{equation}\label{jy1}
    C(X)=\{0\},
\end{equation}
or
\begin{equation}\label{jy2}
    C(X)=\{0,\operatorname{diam}X\}
\end{equation}
and, moreover, the last equality holds if and only if 
\begin{equation}
    \label{er0}
\operatorname{diam} X \in D(X).
\end{equation}



%Equality \eqref{jy2} holds if and only if $\operatorname{diam}X\in D(X)$.

\end{Theorem}

\begin{Proposition}\label{sat}
Let $(X,d)$ be a metric space. Then the singleton $\{c\}$ is a centered sphere in $(X,d)$ for every $c\in X$.
\end{Proposition}

\begin{Theorem}\label{shak1}
Let $(X,d)\in {\bf  UT}$ have at least two points and let $C(X)$ be the center of distances of $(X,d)$.
Then the equality
\begin{equation}
    \label{ks1-1}
C(X)=\{0,\operatorname{diam}X\}
\end{equation}
holds if and only if $X$ is a centered sphere in $(X,d)$ and $\operatorname{diam} X$ is the radius of this sphere.
\end{Theorem}

\begin{Lemma}\label{scgkk}
Let $(X,d)\in \mathbf{UT}$ and let $T=T(l)$ be a labeled tree such that $(X,d)$ and $(V(T),d_l)$ are isometric.
Then, for every $B^{1}\in {\bf B}_{X}$ there is a subtree $T^{1}$ of $T$ such that $
(B^{1},d\big|_{B^{1}\times B^{1}})$
and
$(V(T^1), d_{l^{1}})$ are isometric,
 where $l^{1}$ is the restriction of the labeling
$l : V(T) \to \mathbb{R}^{+}$ on the set $V(T^{1})$.
\end{Lemma}

\begin{Lemma}\label{dhiu}
     Let $(X,d)$ be an ultrametric space and let $B$ be an open ball in $(X,d)$. Then $B$ is a centered sphere in $(X,d)$ if and only if this ball is a centered sphere in $(B, d|_{B\times B})$.
\end{Lemma}

\begin{Theorem}\label{efhik}
Let $(X,d)$ be an ultrametric space and let 
\begin{equation}
    \label{fesdi6-1}
D_p(X) := \{ d(p,x) : x \in X \}
\end{equation}
for each $p\in X$.
If the inclusion 
\begin{equation}
    \label{ed1}
{\bf B}_X \subseteq {\bf Cs}_X
\end{equation}
holds, then the set
$D_p(X)\cap [0,r)$ has the greatest element for all $p\in X$ and $r\in (0,\infty)$. Moreover, if $(X,d)\in {\bf UT}$ and the set $D_p(X)\cap [0,r)$ has the greatest element for all $p\in X$ and $r\in (0,\infty)$, then inclusion \eqref{ed1} holds.
\end{Theorem}

\begin{proof}
Let inclusion \eqref{ed1} hold, let $p$ be an arbitrary point of $X$, and let
$r\in(0,\infty)$.
We must show that the set $D_p(X)\cap[0,r)$ has the greatest element. 

Let us consider the open ball
\begin{equation}\label{oo1}
    B:=\{x\in X \colon d(x,p)<r\}.
\end{equation}
 By inclusion \eqref{ed1} the ball $B$ is a centered sphere in $(X,d)$, i.e., the equality
\[
B=\{x\in X\colon \ d(x,p_1)=r_1\}\cup\{p_1\}
\]
holds with some $p_1\in B$ and $r_1\in \mathbb R^+$. The last equality and the Lemma~\ref{nnhh}  with $A=B$ and $p=p_1$ imply 
\[
\operatorname{diam} B=r_1 \quad \text{and} \quad
 \operatorname{diam} B\in D(B),
\]
where $D(B)$ is the distance set of $(B,d|_{B\times B})$.

Consequently, the diametrical graph \(G_B\) of \((B,d|_{B\times B})\) is non-empty.
Theorem \ref{ref} implies that \(G_B\) is a complete multipartite graph.
Hence, by Definition \ref{scfr}, there is a point \(x\in B\) such that
\begin{equation*}
d(p,x)=\operatorname{diam} B.
\end{equation*}
Using \eqref{fesdi6-1}, \eqref{oo1} and formula \eqref{lokias} with \(A=B\), we see that \(\operatorname{diam} B\) is the greatest element of the set \(D_p(X)\cap[0,r)\) as required.


Let $(X,d)\in {\bf UT}$ and let the set $D_p(X)\cap [0,r)$ contain the greatest element for all $p\in X$ and $r\in (0,\infty)$.
Let us consider an arbitrary open ball $B^1\in {\bf B}_X$ with  a center $p_1\in X$ and a radius $r_1>0$,
\begin{equation}
    \label{nn1}
    B^1=\{x\in X \colon d(p_1,x)<r_1\}.
\end{equation}
To prove inclusion \eqref{ed1} we must show that $B^1$ is a centered sphere in $(X,d)$.

If $B^1$ is a singleton, then $B^1\in {\bf Cs}_X$ holds by Proposition~\ref{sat}. 

Suppose now that $B^1$ contains at least two points.
Since $(X,d)$ is an ${\bf UT}$-space, Lemma~\ref{scgkk} implies the existence of a labeled tree
$T^{1}=T^{1}(l_{1})$ such that
$(B^1,d|_{B^{1}\times B^{1}})$ and $(V(T^{1}),d_{l_1})$
are isometric.
Hence the ultrametric space $(B^{1},d|_{B^{1}\times B^{1}})$ also belongs to the class
${\bf UT}$.

By Theorem~\ref{shak1}, the membership
\((B^{1},d|_{B^{1}\times B^{1}})\in {\bf UT}
\)
implies  \(B^{1}\in  {\bf Cs}_{B^1}\)  if and only if
\begin{equation}
    \label{ng1}
C(B^{1})=\{0,\operatorname{diam} B^{1}\},
\end{equation}
where $C(B^1)$ is the center of distances of $(B^1, d|_{B^1\times B^1})$.
By Theorem~\ref{dggkq}, equality \eqref{ng1} is equivalent to 
\begin{equation}
\label{ng2}
\operatorname{diam} B^{1} \in D(B^{1}),
\end{equation}
where \(D(B^{1})\) is the distance set of the space \((B^{1},d|_{B^{1}\times B^{1}})\).
The space \((B^{1},d|_{B^{1}\times B^{1}})\) is a subspace of \((X,d)\).
Hence the inclusion
\[
D(B^{1}) \subseteq D(X)
\]
holds. Let $D_{p_1}(X)$ be the set defined by \eqref{fesdi6-1} with $p=p_1$, where $p_1$ is the center of the ball $B^1$. By supposition the set $D_{p_1}(X)\cap [0,r_1)$  has the greatest element $R_1$. Thus we have 
\begin{equation*}
    R_1\in D_{p_1} (X)\cap [0,r_1) \subseteq D(B^1).
\end{equation*}
Lemma~\ref{nnhh} and \eqref{nn1} imply $R_1=\operatorname{diam}B^1$.
  Thus \eqref{ng2} holds and, consequently, we have
\begin{equation}
    \label{ng3}
 B^1\in {\bf Cs}_{B^{1}},
\end{equation}
by Theorems \ref{dggkq} and \ref{shak1}.
Using Lemma \ref{dhiu} we may rewrite \eqref{ng3} as
\begin{equation}
    \label{ng4}
B^1 \in {\bf Cs}_X. 
\end{equation}
Thus \eqref{ng4} holds for each \(B^{1}\in {\bf B}_X\).
Inclusion \eqref{ed1} follows.

The proof is completed.
\end{proof}

\end{document}
